\section{Future Directions}
\label{sec:future_directions}

The contributions of this dissertation open several avenues for future research. These directions concern the scalability and integration of structured abstractions, the formal and evaluative foundations of continuing adaptation, and the extension of the developed methods to perceptual, multi-agent, and lifelong settings.



\paragraph{Scaling structured abstractions to richer domains.}
The benchmark environments and adaptation methods developed in this dissertation were evaluated in gridworld and game-based domains that, while supporting diverse classes of novelty, operate at a level of complexity below real-world deployment settings. Recent neuro-symbolic methods have begun to jointly acquire continuous skills and symbolic abstractions from limited demonstrations and from failure-recovery experience~\citep{lorang2025fewshot,li2026recover}. These methods provide concrete mechanisms for extending hybrid planning and learning to continuous state spaces, partial observability, and richer task structures. The force-space methods in \cref{ch:flex,ch:pho} operate in continuous robotic manipulation settings over restricted task classes involving articulated-object manipulation and planar pushing. A natural integration direction is to embed force-based skills within multi-step manipulation plans, allowing the action abstractions developed for adaptation and the interaction abstractions developed for physical generalization to operate within a single agent architecture.

\paragraph{Integrating perception with adaptation and generalization.}
The adaptation framework in \cref{ch:rapid_learn} largely assumes access to task-relevant symbolic state, while the robotic methods in \cref{ch:flex,ch:pho} use task-specific perception and physical probing to estimate the information required for execution. Recent neuro-symbolic approaches learn symbol grounding from planning constraints and interaction, train perceptual components online from instruction, and invent visually grounded predicates for black-box skills~\citep{lymperopoulos2025acttoground,schneider2025instruction,yang2025skillwrapper}. Integrating these capabilities with novelty detection, localization, and recovery would enable the agent to revise its perceptual representations as part of adaptation. Future work could relate perceptual mismatch---including previously unseen objects, attributes, relations, and sensing conditions---to the task-impact categories in \cref{sec:novelty_taxonomy} and determine whether recovery requires changes to perception, the symbolic model, an executor, or several components.

\paragraph{Bridging structured abstractions and foundation models.}
\Cref{sec:foundation_models} identified a complementary relationship between structured abstractions and foundation models. Recent systems use vision-language models to generate constraints for task and motion planning, large language models to propose missing operators and reward functions for learning, and foundation models to invent grounded predicates for robot skills~\citep{kumar2026openworldtamp,lu2026novelty,yang2025skillwrapper}. Future research can incorporate these capabilities into a verified adaptation loop in which generated representations are evaluated through planning and physical interaction, revised when execution evidence contradicts them, and retained when they support reliable transfer. Generalist vision-language-action policies~\citep{team2024octo,black2024pi_0} could supply initial executors within the operator--executor framework, with the targeted learning mechanisms of \cref{ch:rapid_learn} refining them after localized failures. Force-space representations could provide a corresponding interface for contact-rich execution by connecting generated task knowledge to measurable physical interaction dynamics.

\paragraph{Formal connections between novelty type and recovery strategy.}
The novelty taxonomy of \cref{sec:novelty_taxonomy} classifies novelties by their impact on task performance (prohibitive, obstructive, beneficial, nuisance), and the adaptation framework of \cref{ch:rapid_learn} provides a cascading recovery pipeline. Recovery selection could be formulated as a decision problem over diagnostic uncertainty, expected adaptation cost, likelihood of recovery, safety constraints, and the consequences of unsuccessful exploration. Goal reasoning and information-directed exploration provide mechanisms for reasoning about impasses and selecting informative experience~\citep{aha2018goalreasoning,lorang2022speeding}, while recent work on safe continual reinforcement learning demonstrates the difficulty of preserving safety and previously acquired competence during adaptation to nonstationary dynamics~\citep{coursey2026safe}. Future theory could establish sufficient conditions under which replanning, model revision, executor adaptation, or exploratory action is appropriate for a diagnosed novelty and could characterize how these conditions change as evidence accumulates.

\paragraph{Multi-agent and collaborative open-world settings.}
The framework developed in this dissertation considers a single agent operating in an environment. Multi-agent open-world settings introduce agent-specific observations of novelty, distributed model revision, communication constraints, and interactions among concurrently adapting policies. The benchmark environments could support controlled studies of asymmetric novelty exposure, shared recovery knowledge, and coordination failure. In physical manipulation, existing work on collaborative planar pushing demonstrates the joint planning of contact modes and forces for teams of robots~\citep{tang2024collaborative}. The force-space representations in \cref{ch:flex,ch:pho} provide a compatible basis for studying how multiple robots coordinate interaction forces, revise their roles after an unexpected change, and share evidence about altered object dynamics.

\paragraph{Lifelong learning and knowledge accumulation.}
The adaptation methods in \cref{ch:rapid_learn} support recovery from individual novelty events and the reuse of learned executors. Lifelong operation additionally requires the agent to accumulate, organize, and revise knowledge across a continuing sequence of novelty encounters. Neurosymbolic continual learning and lifelong creative problem solving provide initial mechanisms for retaining and reusing structured knowledge across changing tasks~\citep{lorang2024goalconditioned,gizzi2022lifelong}. Recent work also formulates continual reinforcement learning through history processes and deviation regret, emphasizing the evaluation of evolving behavior during non-resettable interaction~\citep{elelimy2025rethinking}. These developments motivate evaluation protocols in which novelties arise sequentially and the agent's knowledge, behavior, and resource use are evaluated throughout its deployment history. Within the structured abstraction framework, important research problems include deciding which models and executors to retain, identifying when prior adaptations can be transferred or refined, and limiting interference among accumulated capabilities. Human guidance could further support knowledge alignment by helping the agent relate newly acquired representations to existing concepts and determine when previous adaptations are relevant to a novel scenario.
